On 20 November 2026 Levent Alp\"oge released a Lean~4 formalisation, \linemillions{} million
lines long, of a proof of the Riemann Hypothesis and of its analogue for every Dirichlet
$L$-function, with a note saying that no human-readable proof has been written. We built the
archive, twice. It compiles, the five top theorems depend on no axiom beyond the three on
which Mathlib rests, and the kernel accepts every theorem, and every definition that is not
marked unsafe or partial, of the import closure, \replaymillions{} million declarations in
all, replayed from an empty environment. This
paper says what the argument is, as read from the Lean source. The exponent $\tfrac12$ comes
from one optimisation: an upper bound, of the shape $x/m+m$, for a weighted count of prime
powers whose Frobenius lies outside a CM type, where $m$ is the length of a chain of theta
lifts between definite unitary groups. The bound is a squeeze between two slopes of Hermitian
lattices. Divisibility of periods at inert primes, which comes from the height-two formal group
of a CM elliptic curve, raises the slopes of lattices of algebraic modular forms. A slope
inequality for nearly tight frames of maps, which we state in full, transfers that to a lattice
of periods. The slope of that lattice is bounded above by a period computation that returns
exactly the main term of the count, through a dimension count and the Chowla--Selberg formula. We describe each step
and each input, determine by machine which parts of the archive the final theorem uses, and
identify the published theorems on which the proof rests.
